\section{Full Prompt Templates}
\label{app:prompts}

\DIFaddbegin
This appendix lists the templates used by the main LLM call sites referenced in Section~\ref{sec:prompt}, abridged to their static skeleton and runtime slots (the verbose in-prompt admonitions are elided). Each template's static body is shown in black; runtime-injected slots appear in gray with colored tags from the shared figure palette: \DiffTag{figloop}{LOOP/SYMBOLIC} for loop-aware symbolic context, \DiffTag{figcall}{ROUTE/CALL} for route and call-site context, and \DiffTag{figcontract}{VERIFIER/SPECIFICATION} for verifier errors and specification edits.
\DIFaddend


\DIFbluePromptSubsection{Loop-invariant generation}{app:prompt-loop}

\DIFaddbegin
The example below is a filled loop-invariant prompt for a small branch-then-loop program, with the long example and guidance blocks elided.
\DIFaddend

\begin{tcolorbox}[promptboxcompact, title={Loop-invariant generation prompt}]
\begin{lstlisting}[style=PromptInnerCompact]
### Role / Task ###
Given a C program with a loop, generate ACSL loop invariants
that let Frama-C verify it.

(*@\DiffTag{figcall}{ROUTE}\quad{\normalfont\itshape\color{gray!60!black}examples + ACSL helper policy for the matched route, omitted}@*)
(*@\DiffTag{gray}{GUIDANCE}\quad{\normalfont\itshape\color{gray!60!black}strength / supplementary / assertion guides, omitted}@*)

### Rules ###
- Only fill the loop-invariant PLACE_HOLDERs (or add a helper block
  before the function); keep existing annotations unchanged.
- Use the pre-condition, not \at(var, LoopEntry).
- Emit one self-contained ```c``` block.

(*@\DiffTag{figloop}{SYMBOLIC STATE}\quad{\normalfont\itshape\color{gray!60!black}supplied by the symbolic executor}@*)
Pre-condition:
    (\at(a,Pre)==0 && y==0 && x==1 && a==\at(a,Pre))
 || (\at(a,Pre)!=0 && y==0 && x==0 && a==\at(a,Pre))

Loop + skeleton to complete:
/*@
  loop invariant (\at(a,Pre)!=0) ==> ((y<100000) ==> PLACE_HOLDER_x);
  loop invariant (\at(a,Pre)!=0) ==>
                 (!(y<100000) ==> (y==0 && x==0 && a==\at(a,Pre)));
  loop invariant (\at(a,Pre)!=0) ==> (a==\at(a,Pre));
  ... (symmetric a==0 branch omitted) ...
*/
while (y < 100000) { x = x + y; y = y + 1; }
\end{lstlisting}
\end{tcolorbox}

\DIFbluePromptSubsection{Precondition fallback}{app:prompt-pre}

\begin{tcolorbox}[promptboxcompact, title={Precondition fallback prompt}]
\begin{lstlisting}[style=PromptInnerCompact]
### Role: ###
Generate ONLY function-level preconditions plus supporting
predicate/logic declarations needed by them.

### Goal: ###
Rule out Frama-C runtime errors:
- pointer dereferences must be covered by \valid or
  \valid_read;
- array ranges must cover all indexed accesses;
- inductive data structures use structural predicates from
  examples.

(*@\DiffTag{figcall}{ROUTE EXAMPLES}\color{gray!65!black}@*)
{examples}

### Output format: ###
Output exactly one fenced ```acsl``` block:
/*@ predicate <name>(...) = ...; */
/*@ logic <type> <name>(...) = ...; */
/*@
  requires <clause>;
*/

### Hard rules: ###
- Output only ACSL annotations.
- Produce requires only: no ensures, assigns, or loop invariant.
- Each predicate or logic declaration lives in its own /*@ ... */ block.

Function under specification:
{code}
\end{lstlisting}
\end{tcolorbox}

\DIFbluePromptSubsection{Postcondition and assignment generation}{app:prompt-post}

\begin{tcolorbox}[promptboxcompact, title={Postcondition / assignment generation prompt}]
\begin{lstlisting}[style=PromptInnerCompact]
### Task: ###
Transform the following C code into an ACSL-annotated version.

### Rules: ###
- Only fill ensures PLACE_HOLDER and assigns PLACE_HOLDER.
- The input is a complete C file; callee specifications are
  read-only context.
- Output only the modified target function block. The
  harness splices it back into the original file.
- Helper predicate/logic definitions are allowed and required
  if the target specification references a name not already defined.
- Replace the ensures PLACE_HOLDER with as many strong, non-redundant
  clauses as the behavior admits (typically 3-8): the closed-form
  result, per-branch facts, and frame consequences.
- Inline relations directly in ensures; add a predicate/logic helper
  only when the relation is recursive or higher-order.

(*@\DiffTag{figloop}{LOOP-AWARE DERIVATION}\color{gray!65!black}@*)
{loop_guide}

### ACSL pitfalls ###
(*@\DiffTag{figcall}{UNIVERSAL + SYMBOLIC-ROUTE HINTS}\color{gray!65!black}@*)
{category_hints}

(*@\DiffTag{figloop}{RESOLVED ARRAY SIZES}\color{gray!65!black}@*)
{array_length_facts}

### C Program Code: ###
{code}
\end{lstlisting}
\end{tcolorbox}

\DIFbluePromptSubsection{Syntax repair}{app:prompt-syntax}

\DIFaddbegin
A deterministic ACSL fixer (Section~\ref{sec:refinement}) runs before this prompt and clears the common mechanical faults; the template below is the LLM fallback used only when the fixer cannot resolve the parser error.
\DIFaddend

\begin{tcolorbox}[promptboxcompact, title={Syntax-repair prompt}]
\begin{lstlisting}[style=PromptInnerCompact]
### Role: ###
You are an expert in C and Frama-C.

### Task: ###
Correct syntactically incorrect ACSL annotations using
Frama-C errors.

Syntax error:
(*@\DiffTag{figcontract}{VERIFIER ERROR}\color{gray!65!black}@*)
{error_str}

C code with incorrect ACSL annotations:
{c_code}

### ACSL pitfalls ###
(*@\DiffTag{figcall}{UNIVERSAL + SYMBOLIC-ROUTE HINTS}\color{gray!65!black}@*)
{category_hints}

### Rules: ###
- Do not modify function bodies or signatures.
- Only change ACSL annotations.
- New top-level predicate/logic blocks may be added before
  the function.
- If an unbound logic/predicate name is reported, define it
  or remove every clause that references it.
\end{lstlisting}
\end{tcolorbox}

\DIFbluesubsection{Pairwise specification strength judge}
\label{app:prompt-strength-judge}

\DIFaddbegin
Fig.~\ref{fig:prompt-judge} gives the prompt used for the pairwise strength comparison reported in Table~\ref{tab:strength_summary} (Section~\ref{subsec:rq6-strength}).
\DIFaddend

\begin{figure}[!htbp]
\centering
\begin{tcolorbox}[promptboxcompact, title={Pairwise strength-judge prompt}, width=\columnwidth]
\begin{lstlisting}[style=PromptInnerCompact]
[System]
You compare two formal-specification artifacts for the same C function.
Both have already been accepted by their respective verifier. Judge which
one is the STRONGER (more informative) functional specification.

[User]
The same function has been annotated by two tools. Both specifications have been
verified by their respective verifier. Your task is to compare the STRENGTH
of the specifications: which one more precisely characterizes the function's
input-output behavior?

Strength rubric (apply in order, the first applicable rule decides):
  1. STRONGER  = if A's ensures/loop-invariants logically imply B's, AND A's
                 contain at least one non-trivial fact that B's do not.
  2. WEAKER    = symmetric to STRONGER.
  3. EQUAL     = both specifications express the same set of facts (modulo
                 alpha-renaming, syntactic reordering, or trivially redundant
                 clauses).
  4. INCOMPARABLE = each specification has at least one fact the other lacks, and
                    neither chain of implication holds.

Ignore: pretty-printing, comments, syntax style, frame-condition formatting
(\nothing vs. explicit \assigns lists are equal here), and which variables
are mentioned via \at(.,Pre) versus parameters.

Heavily favor postcondition/loop-invariant content over precondition content,
because the precondition surface is set by the verifier configuration.

Reply with strict JSON, no markdown, no commentary, schema:
{"verdict": "A_stronger" | "B_stronger" | "equal" | "incomparable",
 "reason": "<one sentence>"}

A = {tool_a}
B = SESpec

A ({tool_a} spec, language: {lang_a}):
```{lang_a}
{spec_a}
```

B (SESpec spec, language: C/ACSL):
```c
{spec_b}
```
\end{lstlisting}
\end{tcolorbox}
\caption{Pairwise strength judge prompt for Table~\ref{tab:strength_summary}. \(A\) is \autospec\ or \specgen, and \(B\) is \toolname.}
\label{fig:prompt-judge}
\end{figure}

\DIFbluesubsection{Verifier-guided semantic refinement}
\label{app:prompt-refine}

\DIFaddbegin
Fig.~\ref{fig:prompt-refine} gives the refinement prompt used by the verifier-guided refinement loop of Section~\ref{sec:refinement}.
\DIFaddend

\begin{figure}[!htbp]
\centering
\begin{tcolorbox}[promptboxcompact, title={Verifier-guided refinement prompt}, width=\columnwidth]
\begin{lstlisting}[style=PromptInnerCompact]
### Task: ###
Frama-C reports verification failures grouped by category. Each category
has a different fix strategy; only relevant strategies are shown.

### Categorized errors ###
(*@\DiffTag{figcontract}{ACTIVE ERROR BUCKETS}\color{gray!65!black}@*)
{error_str}

### Per-category fix strategy ###
(*@\DiffTag{figcontract}{ACTIVE STRATEGIES ONLY}\color{gray!65!black}@*)
{strategy_blocks}

### Current C code with annotations ###
{c_code}

(*@\DiffTag{figloop}{LOOP-AWARE DERIVATION}\color{gray!65!black}@*)
{loop_guide}

### ACSL pitfalls ###
(*@\DiffTag{figcall}{UNIVERSAL + SYMBOLIC-ROUTE HINTS}\color{gray!65!black}@*)
{category_hints}

### Output rules ###
- The target function's requires may be strengthened when necessary.
(*@\DiffTag{figcontract}{EDIT MODE}\color{gray!65!black}@*)
{output_format_block}
- Do not modify user assert statements.
- Do not modify correct loop invariants or correct postconditions.
- Wrap the corrected target block, or the complete file in full-file mode,
  in one ```c``` fenced block.
\end{lstlisting}
\end{tcolorbox}
\caption{Verifier-guided semantic refinement prompt.}
\label{fig:prompt-refine}
\end{figure}
